\textbf{Phần 1. Trắc nghiệm lựa chọn}
\setcounterpageref{ex}{0}

\begin{ex} %[3H3N18]
Dụng cụ nào sau đây thường dùng để kiểm tra và vẽ góc vuông?
\choice
{Com-pa}
{\True Ê-ke}
{Thước cuộn}
{Thước kẹp}
\loigiai{
Thước ê-ke dùng để kiểm tra góc vuông.\\
\textbf{Đáp án của câu hỏi là B}
}
\end{ex}

\begin{ex} %[3H3N18]
Góc gồm có
\choice
{\True 1 đỉnh và 2 cạnh}
{1 đỉnh và 1 cạnh}
{2 đỉnh và 1 cạnh}
{2 đỉnh và 2 cạnh}
\loigiai{
Góc gồm một đỉnh và hai cạnh xuất phát từ đỉnh đó.\\
\textbf{Đáp án của câu hỏi là A}
}
\end{ex}

\begin{ex} %[3H3H18]
Một hình vuông có bao nhiêu góc vuông?
\choice
{$2$ góc vuông}
{$3$ góc vuông}
{$1$ góc vuông}
{\True $4$ góc vuông}
\loigiai{
Hình vuông có 4 góc vuông.\\
\textbf{Đáp án của câu hỏi là D}
}
\end{ex}

\begin{ex} %[3H3H18]
Một hình chữ nhật có bao nhiêu góc vuông?
\choice
{\True $4$ góc vuông}
{$2$ góc vuông}
{$3$ góc vuông}
{Không có góc vuông nào}
\loigiai{
Hình chữ nhật có 4 góc vuông.\\
\textbf{Đáp án của câu hỏi là A}
}
\end{ex}

\begin{ex} %[3H3H18]
Góc tạo bởi kim giờ và kim phút lúc 3 giờ đúng là
\choice
{Góc không vuông}
{\True Góc vuông}
{Góc bẹt}
{Góc tù}
\loigiai{
Lúc 3 giờ đúng hai kim tạo thành một góc vuông.\\
\textbf{Đáp án của câu hỏi là B}
}
\end{ex}

\begin{ex} %[3H3H18]
Góc tạo bởi kim giờ và kim phút lúc 9 giờ đúng là
\choice
{Góc không vuông}
{Góc bẹt}
{\True Góc vuông}
{Không xác định}
\loigiai{
Lúc 9 giờ đúng tạo thành góc vuông.\\
\textbf{Đáp án của câu hỏi là C}
}
\end{ex}

\begin{ex} %[3H3V18]
Một hình tam giác có thể có nhiều nhất bao nhiêu góc vuông?
\choice
{$2$ góc vuông}
{$3$ góc vuông}
{$4$ góc vuông}
{\True $1$ góc vuông}
\loigiai{
Một hình tam giác chỉ có thể có nhiều nhất 1 góc vuông (tam giác vuông).\\
\textbf{Đáp án của câu hỏi là D}
}
\end{ex}

\begin{ex} %[3H3V18]
Lúc 6 giờ đúng, góc tạo bởi kim giờ và kim phút là góc gì?
\choice
{\True Góc không vuông (góc bẹt)}
{Góc vuông}
{Góc nhọn}
{Không có góc}
\loigiai{
Lúc 6 giờ hai kim nằm trên một đường thẳng (góc không vuông, thẳng hàng).\\
\textbf{Đáp án của câu hỏi là A}
}
\end{ex}

\begin{ex} %[3H3V18]
Trong một hình chữ nhật, số góc vuông là
\choice
{$2$}
{\True $4$}
{$6$}
{$8$}
\loigiai{
Hình chữ nhật có 4 góc vuông.\\
\textbf{Đáp án của câu hỏi là B}
}
\end{ex}

\begin{ex} %[3H3V18]
Hình tứ giác có 2 góc vuông và 2 góc không vuông là hình gì sau đây?
\choice
{Hình vuông}
{Hình chữ nhật}
{\True Hình thang vuông (hoặc tứ giác có 2 góc vuông)}
{Hình tam giác}
\loigiai{
Hình thang vuông có 2 góc vuông và 2 góc không vuông.\\
\textbf{Đáp án của câu hỏi là C}
}
\end{ex}

\begin{ex} %[3H3V18]
Khi gấp đôi tờ giấy hình vuông theo đường chéo, ta được hình tam giác có bao nhiêu góc vuông?
\choice
{$0$ góc vuông}
{$2$ góc vuông}
{$3$ góc vuông}
{\True $1$ góc vuông}
\loigiai{
Được 1 hình tam giác vuông có đúng 1 góc vuông.\\
\textbf{Đáp án của câu hỏi là D}
}
\end{ex}

\begin{ex} %[3H3C18]
Một hình gồm 2 hình vuông ghép liền kề nhau tạo thành một hình chữ nhật lớn. Tổng số góc vuông của cả 2 hình vuông nhỏ và hình chữ nhật lớn tạo thành là
\choice
{$8$ góc vuông}
{$10$ góc vuông}
{\True $12$ góc vuông}
{$6$ góc vuông}
\loigiai{
Mỗi hình vuông nhỏ có 4 góc vuông ($4 \times 2 = 8$). Hình chữ nhật lớn ngoài cùng có 4 góc vuông $\Rightarrow$ Có tổng cộng 12 góc vuông (tính cả các góc thành phần).\\
\textbf{Đáp án của câu hỏi là C}
}
\end{ex}

\textbf{Phần 2. Trắc nghiệm đúng sai}
\setcounterpageref{ex}{0}

\begin{ex} %[3H3H18]
Xét tính đúng sai về các góc:
\choiceTF
{\True Ê-ke có một góc vuông và hai góc không vuông.}
{\True Khung cửa sổ hình chữ nhật có 4 góc vuông ở 4 góc.}
{Lúc 1 giờ đúng, hai kim đồng hồ tạo thành góc vuông.}
{\True Lúc 1 giờ đúng, hai kim đồng hồ tạo thành góc không vuông.}
\loigiai{
\begin{itemize}
    \item Ê-ke có 1 góc vuông: Đúng.
    \item Khung cửa sổ chữ nhật có 4 góc vuông: Đúng.
    \item Lúc 1 giờ là góc nhọn (không vuông). Ý 3 Sai, Ý 4 Đúng.
\end{itemize}
\begin{itemize}
    \item a) \textbf{Đúng}.
    \item b) \textbf{Đúng}.
    \item c) \textbf{Sai}.
    \item d) \textbf{Đúng}.
\end{itemize}
}
\end{ex}

\begin{ex} %[3H3V18]
Cho hình tam giác $ABC$ có góc đỉnh $A$ là góc vuông. Xét tính đúng sai:
\choiceTF
{\True Tam giác $ABC$ có 1 góc vuông.}
{\True Tam giác $ABC$ có 2 góc không vuông (tại đỉnh $B$ và đỉnh $C$).}
{Tam giác $ABC$ có 2 góc vuông.}
{\True Hai cạnh tạo thành góc vuông là $AB$ và $AC$.}
\loigiai{
\begin{itemize}
    \item Có 1 góc vuông: Đúng.
    \item 2 góc còn lại không vuông: Đúng.
    \item Không thể có 2 góc vuông. Sai.
    \item Hai cạnh góc vuông là $AB$ và $AC$: Đúng.
\end{itemize}
\begin{itemize}
    \item a) \textbf{Đúng}.
    \item b) \textbf{Đúng}.
    \item c) \textbf{Sai}.
    \item d) \textbf{Đúng}.
\end{itemize}
}
\end{ex}

\textbf{Phần 3. Trả lời ngắn}
\setcounterpageref{ex}{0}

\begin{ex} %[3H3H18]
Hình vuông có bao nhiêu góc vuông?
\shortans{4}
\loigiai{
\textbf{Đáp số:} $4$\\
4 góc vuông.
}
\end{ex}

\begin{ex} %[3H3H18]
Hình chữ nhật có bao nhiêu góc không vuông?
\shortans{0}
\loigiai{
\textbf{Đáp số:} $0$\\
Cả 4 góc đều là góc vuông nên có 0 góc không vuông.
}
\end{ex}

\begin{ex} %[3H3V18]
Lúc mấy giờ đúng vào buổi chiều thì kim giờ và kim phút tạo thành một góc vuông (chỉ ghi số giờ từ 1 đến 12)?
\shortans{3}
\loigiai{
\textbf{Đáp số:} $3$\\
Lúc 3 giờ (hoặc 9 giờ).
}
\end{ex}

\begin{ex} %[3H3V18]
Một hình tam giác vuông có bao nhiêu góc không vuông?
\shortans{2}
\loigiai{
\textbf{Đáp số:} $2$\\
Có 2 góc không vuông.
}
\end{ex}

\begin{ex} %[3H3V18]
Một hình ngũ giác có 5 góc, trong đó có 2 góc vuông. Hỏi có bao nhiêu góc không vuông?
\shortans{3}
\loigiai{
\textbf{Đáp số:} $3$\\
$5 - 2 = 3$ góc không vuông.
}
\end{ex}

\begin{ex} %[3H3C18]
Có 3 hình vuông xếp cạnh nhau. Tổng số góc vuông của cả 3 hình vuông đó là bao nhiêu?
\shortans{12}
\loigiai{
\textbf{Đáp số:} $12$\\
$4 \times 3 = 12$ góc vuông.
}
\end{ex}
